\begin{table}[htbp]
\caption{Treatment-Control Differences in Selected Covariates and Sample Indicators: Liberia}
\label{tab:lib_balance}
\begin{center}
\begin{small}
\begin{threeparttable}
\begin{tabular}{l*{5}{r}}
\toprule
 & \Longstack{Treatment\\ mean} & \Longstack{Control\\ mean} & Difference & \Longstack{Standard\\ error} & $p$-value \\
\midrule
Baseline score & 19.848 & 21.014 & -1.134 & (1.008) & 0.260 \\
Upper group & 0.549 & 0.612 & -0.056 & (0.045) & 0.216 \\
Has endline & 0.660 & 0.659 & 0.002 & (0.062) & 0.976 \\
Class size & 59.278 & 45.596 & 13.519 & (4.331) & 0.002 \\
\bottomrule
\end{tabular}
\begin{tablenotes}
\footnotesize \item Notes: The unit of observation is a student in the Liberia analytic sample. The table reports treatment and control means together with the treatment-control difference from a student-level OLS regression that includes strata fixed effects. Standard errors are cluster-robust at the school-grade-group level. `Upper group' is the share of students whose baseline score lies above the study cutoff for the relevant grade block, `Has endline' is the share with a matched endline outcome, and `Class size' is realized reading-class size. The table is descriptive and is used to assess whether treatment and control samples differ systematically on key observable margins.
\end{tablenotes}
\end{threeparttable}
\end{small}
\end{center}
\end{table}